% !TeX root = proyecto.tex

%=========================================================
%\chapter{Project Charter}

\newcommand{\ESCOMPchSec}[1]{\rowcolor{colorAgua}\multicolumn{4}{|c|}{\bf #1}\\\hline}
\newcommand{\ESCOMPchItem}[2]{{\bf {#1}} & \multicolumn{3}{p{.66\textwidth}|}{#2}\\\hline}
\newcommand{\ESCOMPchSubItem}[3]{{\bf {#1}} & {#2} & \multicolumn{2}{p{.44\textwidth}|}{#3}\\\hline}
\newcommand{\ESCOMPchSubSubItem}[4]{{\bf {#1}} & {#2} & {#3}& {#4}\\\hline}

\cleardoublepage
{\centering{\Huge Project Charter}\bigskip\\}
\begin{table}[hptb!] 
%\renewcommand\thetable{i}
\begin{tabular}{|p{.22\textwidth} |p{.22\textwidth} |p{.22\textwidth} |p{.22\textwidth} |}
	\hline
	\ESCOMPchItem{Proyecto:}{2026-A135: Sistema Inteligente para la Identificación y Seguimiento de NNA}
	\ESCOMPchItem{Responsable:}{Héctor Morales, Desarrollador Principal}
	\ESCOMPchItem{Autoriza:}{Comisión de Trabajos Terminales (Academia)}
	\ESCOMPchItem{Background/Contexto:}{Miles de NNA quedan en orfandad por feminicidios anualmente; actualmente no existe un sistema centralizado de seguimiento para apoyarlos.}
	\ESCOMPchItem{Beneficios esperados:}{Visibilidad del fenómeno, reducción de trabajo manual para ONGs y generación de base de datos histórica mediante NLP.}
	\ESCOMPchItem{Costo estimado:}{Proyecto Académico (TT1/TT2)}
	\ESCOMPchSubSubItem{Fecha de inicio:}{Septiembre 2025}{\bf Fecha de término:}{Junio 2026}
	\ESCOMPchItem{Objetivo:}{Sistema automatizado de recolección, análisis NLP y visualización de noticias sobre NNA afectados por feminicidios.}
	\ESCOMPchSec{Entregables Principales}
	\ESCOMPchSubItem{}{SIST-01}{Sistema web y pipeline NLP funcional (Contenedores Docker)}
	\ESCOMPchSubItem{}{DOC-01}{Documentación Técnica Final (este documento)}
	\ESCOMPchSubItem{}{COD-01}{Código fuente completo en repositorio}
	\ESCOMPchSec{Alcance del proyecto}
	\ESCOMPchItem{Incluye:}{
		\begin{Titemize}
			\Titem Recolección automatizada multi-fuente (RSS, Google Search).
			\Titem Scraping web dinámico y evasión heurística.
			\Titem Scoring de 4 ejes mediante Zero-Shot y Fine-Tuning.
			\Titem Clasificación semántica con modelo BETO.
			\Titem Agrupación de noticias (clustering) mediante BERTopic.
			\Titem Deduplicación cross-site avanzada.
			\Titem Base de datos PostgreSQL y Dashboard interactivo.
		\end{Titemize}
	}
	\ESCOMPchItem{Excluye:}{
		\begin{Titemize}
			\Titem Identificación y almacenamiento de datos personales reales de NNA.
			\Titem Intervención social directa o canalización oficial a las autoridades.
			\Titem Validación pericial de la veracidad de los hechos periodísticos.
		\end{Titemize}
	}
	\ESCOMPchItem{Criterio de éxito:}{Detección de notas verdaderas con menos del 5\% de falsos positivos.}
	\ESCOMPchItem{Metodología:}{Desarrollo Evolutivo Incremental (Iteraciones de Prototipo P0 a P4)}
	\ESCOMPchSec{Datos de contacto}
	\ESCOMPchItem{Desarrollador:}{Héctor Morales (Autor del TT)}
	\ESCOMPchItem{Directores:}{Directores asignados del TT}
	\ESCOMPchItem{Organización de Apoyo:}{Fundación Futuro con Derechos (Asesoría Social)}
	\ESCOMPchItem{Riesgos y peligros:}{
		\begin{Titemize}
			\Titem Bloqueos anti-scraping progresivos por parte de los medios periodísticos.
			\Titem Lentitud en la inferencia del modelo neuronal por falta de hardware GPU.
		\end{Titemize}
	}
	\ESCOMPchItem{Supuestos:}{
		\begin{Titemize}
			\Titem Los medios noticiosos mantienen un estándar semántico en su redacción de crímenes.
			\Titem Los servicios de búsqueda como Google News permanecen públicamente consultables.
		\end{Titemize}
	}
	\ESCOMPchItem{Restricciones y dependencias:}{
		\begin{Titemize}
			\Titem El sistema no puede operar en tiempo real exacto; opera por lotes (cron-jobs).
			\Titem Restricción absoluta de protección de identidad del menor por normativas vigentes.
		\end{Titemize}
	}
\end{tabular}
	\caption{Resumen del proyecto (Project Charter)}
	\label{tbl:projectCharter}
\end{table}
